\documentclass[10pt,a4paper]{article}
\usepackage[utf8]{inputenc}
\usepackage[T1]{fontenc}
\usepackage[margin=18mm]{geometry}
\usepackage{longtable}
\usepackage{hyperref}
\usepackage{enumitem}
\setlist[itemize]{noitemsep, topsep=2pt}
\setlength{\parindent}{0pt}
\begin{document}

\section*{Non-Verbal \& Multimodal Discourse Analysis Worksheet}

\textit{Qualitative analysis of pauses, trailing sentences, prosody/intonation, vocalisations (coughs, sighs, scoffs), and kinesics (hand gestures, facial expressions, posture) in political news discourse}\par\smallskip

\bigskip\hrule\bigskip

\subsection*{Guidelines: Five non-verbal \& paralinguistic features to look for in every event}

\begin{longtable}{|p{0.23\textwidth}|p{0.23\textwidth}|p{0.23\textwidth}|p{0.23\textwidth}|}
\hline
Feature & What to look for & Illustrative example (from sample) & Questions to answer \\ \hline \endhead
\textbf{1. Pauses, Silences \& Abrupt Stops} & Micro-pauses \texttt{(.)}, timed silences \texttt{(1.5)}, pregnant/dramatic pauses before a loaded word, hesitant stumbles, abrupt self-interruptions or cut-offs (\texttt{---}), and forced silences (muting a panelist's mic or talking over them). & Anchor pauses \texttt{(1.5)} before saying \textit{"'loota'"} for dramatic weight; Addl. DCP cuts off mid-sentence (\textit{"wo toh--- abhi investigation chal raha hai"}) when asked if police struck first; panelist's mic faded during counter-question. & Where does the pause or stop occur? Is it deliberate/rhetorical (for suspense, emphasis, or gravitas) or defensive/evasive (cognitive load, self-repair, avoiding blame)? Who controls turn-taking and silence? \\ \hline
\textbf{2. Trailing Sentences \& Unfinished Utterances} & Clauses that trail off (\texttt{...}), drop in volume at the end, or leave the conclusion unspoken (aposiopesis) so the audience supplies the implication without the speaker taking explicit responsibility. & \textit{"Aur phir wahan jo log ghuse the... aap samajh hi sakte hain..."}; \textit{"500 salon ke sangharsh ke baad agar aisa..."} trailing off with a slow head shake. & What exact thought is left unfinished? What presupposition or shared bias is the viewer invited to fill in? Does trailing off act as a safe hedge to imply higher-level complicity or "anti-national" motive without stating it? \\ \hline
\textbf{3. Intonation, Pitch, Volume \& Prosody} & Pitch contour (rising \texttt{$\uparrow$} for incredulity/challenge; falling \texttt{$\downarrow$} for finality), volume shifts (shouting vs. hushed whisper), emphatic word stress, speech rate (rapid-fire urgency vs. slow solemnity), and tonal register (sarcastic, reverential, combative). & Hushed, reverential cadence on \textit{"Prabhu Shri Ram ke Mandir"} vs. sharp, high-pitched interrogation on \textit{"par Akhilesh ji Ram Mandir kyon$\uparrow$ nahin jate?!"}; satirical sing-song inflection on \textit{"priventiv arenjment"}. & Which exact words receive pitch peaks or volume stress? How does vocal tone turn neutral words into irony, mockery, or moral outrage? Does the speaker shift prosodic register depending on which actor is discussed? \\ \hline
\textbf{4. Coughs, Sighs \& Non-Lexical Vocalisations} & Audible sighs (\texttt{[sigh]}), scoffs/snorts (\texttt{[scoff]}), ironic chuckles (\texttt{[chuckle]}), throat-clearing/coughs (\texttt{[cough]}), sharp intakes of breath/gasps (\texttt{[gasp]}), and tongue-clicks (\texttt{[tsk]}). & Heavy \texttt{[sigh]} before reporting Trust oversight lapses (signalling reluctant sorrow rather than hostility); mocking \texttt{[scoff]} or \texttt{[chuckle]} when reading an opponent's denial; \texttt{[cough / throat-clear]} when pivoting away from an awkward fact. & What emotion or epistemic stance does the vocalisation perform (exasperation, moral disgust, disbelief, ridicule, discomfort)? Does a sigh or scoff pre-judge a speaker's claim before any verbal rebuttal is made? \\ \hline
\textbf{5. Hand Gestures, Facial Expressions \& Posture (Kinesics)} & Hand movements (air-quotes, finger-pointing/wagging, desk-thumping, open palms, counting on fingers, dismissive wave), facial cues (raised eyebrows, eye-rolls, smirk, furrowed brow), head nods/shakes, leaning into camera, and visual props (bandaged limbs, documents). & Anchor forms physical air-quotes on \textit{"'loota'"} or \textit{"'peaceful' protest"}; host points finger/pen at lens during an ultimatum; injured DCP holds bandaged hand up to camera; creator raises one eyebrow and leans in for a conspiratorial reveal. & How do gestures reinforce, contradict, or sharpen the spoken words? During split-screen listening, does the host's face (nodding vs. smirking/eye-rolling) cue the audience on whom to trust? How are bodily visuals used as evidence? \\ \hline
\end{longtable}

\subsection*{Quick Multimodal Transcription Notation Key (use inside transcript excerpts)}

\begin{longtable}{|p{0.15\textwidth}|p{0.15\textwidth}|p{0.15\textwidth}|p{0.15\textwidth}|p{0.15\textwidth}|p{0.15\textwidth}|}
\hline
Symbol & Meaning & Example in transcript & Symbol & Meaning & Example in transcript \\ \hline \endhead
\texttt{(.)} & Brief micro-pause (< 0.5 sec) & \texttt{woh (.) suparavijan lacking tha} & \texttt{$\uparrow$ / $\downarrow$} & Marked rising / falling pitch & \texttt{kyon$\uparrow$ FIR nahin karaee$\downarrow$} \\ \hline
\texttt{(1.5)} & Timed pause in seconds & \texttt{kisne (1.5) 'loota'?} & \texttt{WORD / \_word\_} & Loud volume / emphatic stress & \texttt{ye PAAP hai / \_bilkul\_ galat} \\ \hline
\texttt{---} & Abrupt stop, cut-off, or self-repair & \texttt{police ki taraf se--- abhi jaanch...} & \texttt{$^\circ$word$^\circ$} & Softened / whispered delivery & \texttt{$^\circ$ander ki baat ye hai$^\circ$} \\ \hline
\texttt{...} & Trailing off / unfinished sentence & \texttt{aap jaante hain ki kaun log...} & \texttt{[sigh] [cough] [scoff]} & Non-lexical vocalisations & \texttt{[sigh] 500 salon ke baad...} \\ \hline
\texttt{[overlap]} & Simultaneous / overlapping speech & \texttt{[Anjana: suniye---] [Guest: nahin!]} & \texttt{\{\{gesture: ...\}\}} & Hand gesture / facial cue / posture & \texttt{\{\{air-quotes on 'loota'\}\}} \\ \hline
\end{longtable}

\textbf{Methodological Reminder:} A non-verbal cue is not automatically manipulative or ideological---a cough, pause, or rapid cadence may reflect studio audio delay, live time pressure, or a speaker's habitual style. Always check (1) whether the cue clusters with verbal features, and (2) whether the speaker's delivery shifts noticeably across different actors or events (within-outlet baseline).\par\smallskip

\bigskip\hrule\bigskip

\subsection*{Event 1. Ram Mandir donation controversy (Non-Verbal \& Paralinguistic Layer)}

\begin{itemize}
  \item \textbf{Focus:} How do prosody, pauses, trailing sentences, vocalisations (sighs, coughs, scoffs), and body language frame the controversy---as solemn devotional sorrow, administrative scrutiny, or partisan combat?
  \item \textbf{Outlet:} Dhruv Rathee (YouTube Channel) | \textbf{Date:} 15 June 2021
  \item \textbf{Video title:} Ayodhya Ram Mandir Scam | Exposed by AAP \& SP | Dhruv Rathee | \href{https://youtu.be/oeH0fOSJuX8}{https://youtu.be/oeH0fOSJuX8} | \textbf{Timestamp(s):} \texttt{[00:00]}--\texttt{[15:42]}
  \item \textbf{Analyst:} Baani Solanki (\texttt{2025114016}) | \textbf{Language:} Hinglish (Hindi in Roman script + English terms)
  \item \textbf{Format / Visual setup:} \texttt{[ ]} TV Panel Debate \texttt{[ ]} TV Bulletin \texttt{[x]} Solo YouTube Explainer \texttt{[ ]} Interview / Discussion
  \item \textbf{Primary speaker(s) observed:} Dhruv Rathee (standing/semi-standing studio presenter against a green-screen documentary backdrop, using structured hand-chops, finger-counting, direct camera gaze, and full-screen document highlights)
\end{itemize}

\subsubsection*{Event-specific non-verbal questions (Event 1)}

\noindent\textbf{1.} \textbf{Vocal \& gestural marking of core labels ("scam", "loota", "paap", "gadabdi", "avyavastha"): When the controversy is named, what non-verbal cues accompany the label? (Look for physical air-quotes, vocal scare-quotes, dramatic pauses before the word, hushed/reverential drops in volume, or shouted stress.)}\par\smallskip
\begin{itemize}
  \item \textbf{Notes:}
Unlike TV anchors who place air-quotes around \textit{"scam"} or soften it into \textit{"avyavastha"}, Dhruv Rathee delivers the labels \textit{"scam"}, \textit{"ghotala"}, and \textit{"corruption"} with \textbf{unhedged, assertive falling intonation \texttt{$\downarrow$} and firm downward hand-chops}. In the opening hook (\texttt{[00:18]}), he pauses \texttt{(0.8)} before stressing the phrase: \textit{"Ram Mandir ke chande ke paise mein (0.8) \_bada\_ GHOTALA saamne aaya hai$\downarrow$"} (\texttt{[00:21]}), widening his eyes and bringing both hands down in a sharp, parallel chopping gesture \texttt{\{\{gesture: bilateral vertical hand-chop on 'ghotala'\}\}}. When mentioning the political exploitation of faith (*"bhagwan Ram ke naam par... chanda chori"\texttt{, }[13:45]\texttt{), his voice drops in pitch with a slow, disapproving head shake }\{\{gesture: slow head shake + furrowed brow\}\}`, framing the act not as a minor clerical lapse but as calculated financial fraud.\par\smallskip
\end{itemize}

\noindent\textbf{2.} \textbf{Differential demeanour toward actors (Temple, Trust leadership, counting clerks, donors, opposition accusers): How do the speaker's voice pitch, facial expression, and posture shift when referring to Prabhu Shri Ram / devotees vs. top figures (Champat Rai, Nripendra Misra, CM Yogi) vs. arrested counting staff vs. opposition leaders (Akhilesh Yadav, Awadhesh Prasad)?}\par\smallskip
\begin{itemize}
  \item \textbf{Notes:}
Dhruv Rathee's non-verbal demeanour bifurcates sharply along lines of institutional power versus ordinary citizenry:\par\smallskip
  \item \textbf{Donors / Ordinary Citizens:} Addressed with empathetic warmth, softened volume, and an open palm extended toward the lens (\texttt{"gareeb logon ne apni mehnat ki kamai se sau-sau rupaye diye"}, \texttt{[01:05]}), positioning devotees as earnest victims of deception.
  \item \textbf{Opposition Whistleblowers (Sanjay Singh, Pawan Pandey):} Introduced with level, authoritative courtroom-style cadence (\texttt{[01:42]}), nodding approvingly as their press-conference documents are displayed full-screen.
  \item \textbf{Trust Leadership \& Political Affiliates (Champat Rai, Mayor Rishikesh Upadhyay, Anil Mishra):} Named with a narrowed gaze, raised eyebrows, and pointed index-finger gestures toward the on-screen relationship flowchart (\texttt{[04:15]}--\texttt{[05:30]}). When noting that Trust member Anil Mishra and BJP Mayor Rishikesh Upadhyay acted as official witnesses on BOTH the Rs 2 crore deed and the Rs 18.5 crore agreement five minutes later (\texttt{[04:48]}), Dhruv pauses \texttt{(1.2)}, leans forward, and taps the air twice with his index finger \texttt{\{\{gesture: double index-finger air-tap toward screen\}\}} to signal undeniable collusion.
\end{itemize}

\noindent\textbf{3.} \textbf{Solemn devotional grief vs. high-decibel partisan spectacle: How do pauses, sighs, speech rate, and hand gestures construct the segment's emotional register? Where is delivery slow, hushed, and mournful (sighs, bowed head, folded hands), and where does it escalate into rapid-fire shouting, desk-thumping, or mocking irony?}\par\smallskip
\begin{itemize}
  \item \textbf{Notes:}
Dhruv constructs a \textbf{measured, forensic-prosecutorial register} (\textasciitilde{}135--145 words per minute) that avoids both devotional weeping and studio shouting. His delivery is highly rhythmic and didactic: he speaks in clear, segmented clauses separated by deliberate \texttt{(0.5)} to \texttt{(1.0)} pauses so the viewer can read the highlighted registry PDFs on screen. The emotional register shifts into \textbf{biting, cold irony} during the rebuttal section (\texttt{[08:10]}--\texttt{[11:30]}): rather than shouting, he slows his speech rate down to \textasciitilde{}110 WPM when reading the Trust's excuses, arching his right eyebrow and letting out a quiet, incredulous \texttt{[scoff]} to make the official defence sound self-evidently absurd.\par\smallskip
\end{itemize}

\noindent\textbf{4.} \textbf{Non-verbal evaluation of financial scrutiny: When questions about unreceipted offerings, lack of frisking, or land deals are raised, how does the host/speaker react non-verbally (earnest nodding, leaning into the camera, open-palm appeals vs. eye-rolls, scoffs, dismissive hand waves, or incredulous rising pitch $\uparrow$)?}\par\smallskip
\begin{itemize}
  \item \textbf{Notes:}
Financial scrutiny is the centrepiece of the video and is performed with \textbf{intense pedagogical focus and visual demonstration}:\par\smallskip
  \item When comparing the official circle rate (Rs 5.80 crore) with the actual purchase price (Rs 2 crore) and the Trust's agreement price five minutes later (Rs 18.5 crore, \texttt{[02:50]}--\texttt{[04:10]}), Dhruv uses wide horizontal hand spans \texttt{\{\{gesture: hands held close together for '2 crore', then stretched wide apart for '18.5 crore'\}\}} to spatialise the 825\% markup (\texttt{"ek second mein daam \_nau\_ guna badh gaya$\uparrow$"}, \texttt{[03:38]}).
  \item Conversely, when evaluating the Trust's counter-claim that an earlier unregistered agreement from March 2021 justified the price (\texttt{[09:15]}), he reacts with a skeptical head tilt, an incredulous rising pitch \texttt{$\uparrow$}, and open-palm appeals (\texttt{"toh phir Trust ne seedha Kusum Pathak se zameen kyon$\uparrow$ nahin khareedi?!"}, \texttt{[09:42]}).
\end{itemize}

\noindent\textbf{5.} \textbf{Performing certainty vs. hedging and innuendo (pauses, stops, trailing off): Are allegations delivered with definitive falling intonation $\downarrow$, finger-counting, and pointing at on-screen evidence (CCTV, documents), or hedged with hesitant pauses (.), abrupt self-interruptions (---), throat-clearing [cough], and trailing sentences (...) that imply higher-level complicity without naming it?}\par\smallskip
\begin{itemize}
  \item \textbf{Notes:}
Dhruv performs \textbf{high epistemic certainty} on the documentary trail while using \textbf{strategic trailing questions (\texttt{...})} to implicate higher political protection:\par\smallskip
  \item \textbf{Forensic Certainty:} He counts out the four sequential steps of the land flip on his fingers (\texttt{[02:45]}--\texttt{[04:55]}), pointing directly at the highlighted timestamps (\texttt{7:10 PM} vs. \texttt{7:15 PM}), stamp duty numbers, and witness signatures on screen with definitive falling intonation \texttt{$\downarrow$} (*"Ye sab kuchh sarkari kagazon mein darj hai$\downarrow$"\texttt{, }[05:12]`).
  \item \textbf{Strategic Innuendo via Trailing Off:} When connecting Mayor Rishikesh Upadhyay and witness Anil Mishra to the broader ruling party ecosystem (\texttt{[06:20]}), he slows down and trails off: \textit{"Gawah kaun hain? Trust ke trustee Anil Mishra aur BJP ke Mayor Rishikesh Upadhyay... ab isse zyada saaf saboot kya chahiye ki iske peechhe kaun log..."} (\texttt{[06:32]}), accompanied by a knowing nod \texttt{\{\{gesture: slow knowing nod + raised eyebrows\}\}} that invites the viewer to conclude top-down BJP complicity.
\end{itemize}

\noindent\textbf{6.} \textbf{Turn-taking, interruptions, and acoustic treatment of denials: In debates/interviews: who receives calm, uninterrupted speaking time and sympathetic nods, and who is cut off mid-sentence (---), talked over, or muted? In solo videos: how are clips of official denials framed via reaction pauses, raised eyebrows, or sarcastic chuckles?}\par\smallskip
\begin{itemize}
  \item \textbf{Notes:}
In this solo documentary format, official denials by Champat Rai and Trust spokespersons are never played as unmediated authority; instead, they are \textbf{sandwiched inside a "Myth vs. Reality" prosecutorial frame}. Before presenting each of the Trust's four defences (\texttt{[08:15]}--\texttt{[12:10]}), Dhruv holds up a finger (\textit{"Ab suniye unka pehla bahana"\texttt{, }[08:20]\texttt{) with a slight smirk. After reading Champat Rai's statement that }"hum par toh }100\texttt{ saal se aarop lag rahe hain, Mahatma Gandhi ki hatya ka bhi aarop laga tha"\textit{ (}[11:40]\texttt{), Dhruv freezes for }(1.5)\texttt{ seconds in a deadpan reaction pause }\{\{gesture: 1.5s frozen stare + slow blink\}\}\texttt{, lets out a dry }[scoff]\texttt{, and says with falling disgust: }"Matlab (.) chande ke hisaab ka sawaal poochho (.) toh ye Mahatma Gandhi ki hatya le aate hain$\downarrow$"* (}[11:52]`).\par\smallskip
\end{itemize}

\noindent\textbf{7.} \textbf{Performing the boundary between religious devotion and institutional critique: How does the presenter use vocal tone, sighs, and body language to signal personal reverence for the Temple (or secular objectivity) while criticising---or defending---the Trust's financial management?}\par\smallskip
\begin{itemize}
  \item \textbf{Notes:}
Dhruv performs \textbf{secular constitutional objectivity combined with moral empathy for ordinary believers}. Rather than adopting devotional gestures (folded hands or chanting), he maintains an upright, journalistic posture throughout. However, at \texttt{[13:30]}--\texttt{[14:20]}, when distinguishing between genuine faith (\texttt{shraddha}) and the commercial/political instrumentalisation of religion (\texttt{dhandha}), his voice softens into a earnest, lower-register appeal: \textit{"Jin logon ne shraddha se (.) apni aastha se paise diye..."} (\texttt{[13:38]}), placing an open hand toward the viewer, before hardening his tone and furrowing his brow on \textit{"unko dharam ke naam par \_loota\_ ja raha hai$\downarrow$"} (\texttt{[13:48]}).\par\smallskip
\end{itemize}

\noindent\textbf{8.} \textbf{Self-presentation as the "guardian of public trust": How does the speaker use direct eye-gaze, posture, moral sighs, or emphatic gestures (pointing at the viewer, hands on heart/desk) to position themselves as the authentic protector of devotees or public accountability?}\par\smallskip
\begin{itemize}
  \item \textbf{Notes:}
Dhruv positions himself as the \textbf{fearless independent investigator ("fact-checker") who exposes what mainstream "Godi Media" hides}. At \texttt{[12:30]}--\texttt{[13:15]}, he lets out a heavy moral \texttt{[sigh]}, shakes his head at how TV channels are ignoring the registration documents, and points his index finger directly at the camera lens \texttt{\{\{gesture: direct index-finger point at lens\}\}}: \textit{"Mainstream media aapko ye kagaz \_kabhi\_ nahin dikhayega$\downarrow$ (\texttt{[12:42]}) ... isiliye is video ko har ek insaan tak share kijiye"} (\texttt{[15:10]}). His unblinking eye-gaze and rhythmic pointing construct a direct pact of accountability between creator and citizen.\par\smallskip
\end{itemize}

\subsubsection*{Apply the five non-verbal features (Event 1)}

\begin{longtable}{|p{0.31\textwidth}|p{0.31\textwidth}|p{0.31\textwidth}|}
\hline
Non-verbal feature & Timestamp + spoken anchor words + non-verbal notation & What it does in context (stance, framing, implication) \\ \hline \endhead
\textbf{1. Pauses, silences \& stops} & \texttt{[04:48]} \textit{"Aur sabse badi baat (1.2) dono agreements mein gawah (witness) \_bilkul\_ same hain$\downarrow$"} / \texttt{[11:48]} \texttt{[(1.5) deadpan reaction pause after Champat Rai quote]} & The \texttt{(1.2)} pause before revealing the identical witnesses builds forensic climax; the \texttt{(1.5)} silent stare after Champat Rai's quote frames the Trust's defence as absurd without needing a single shouted word. \\ \hline
\textbf{2. Trailing sentences (\texttt{...})} & \texttt{[06:32]} \textit{"Gawah hain Trustee Anil Mishra aur BJP Mayor Rishikesh Upadhyay... ab aap khud samajh sakte hain ki ye paisa kahan..."} & Leaves the final destination of the Rs 16.5-crore markup unspoken, prompting the viewer to infer organised political kickbacks while staying anchored to the witness names on the deed. \\ \hline
\textbf{3. Intonation, pitch \& volume} & \texttt{[03:38]} \textit{"Sirf \_PANCH\_ minute ke andar (.) zameen ka daam \_NAU\_ guna badh gaya$\uparrow$!"} (sharp rising pitch \texttt{$\uparrow$} + heavy stress on numerals) & Sharp prosodic contrast on \textit{"\_PANCH\_ minute"} vs. \textit{"\_NAU\_ guna"} dramatises the mathematical impossibility of natural price appreciation in 300 seconds. \\ \hline
\textbf{4. Coughs, sighs \& vocalisations} & \texttt{[08:22]} \texttt{[[scoff]]} \textit{"Ab suniye Trust ki taraf se kya safai di gayi hai..."} / \texttt{[12:35]} \texttt{[[sigh]]} \textit{"Par afsos ki baat ye hai ki TV media is par chup hai..."} & The \texttt{[scoff]} pre-frames the Trust's explanation as a flimsy excuse (\textit{"bahana"}), while the \texttt{[sigh]} performs civic sorrow over mainstream media's complicity. \\ \hline
\textbf{5. Hand gestures, face \& posture} & \texttt{[03:35]} \texttt{\{\{gesture: hands close together for '2 crore', then flung wide apart for '18.5 crore'\}\}} / \texttt{[04:50]} \texttt{\{\{gesture: double index-finger tap on witness names\}\}} & Spatialises the financial discrepancy visually for the audience and directs visual attention to the documentary proof (witness signatures) on screen. \\ \hline
\end{longtable}

\subsubsection*{Short summary (Event 1 Non-Verbal Synthesis)}

\begin{itemize}
  \item \textbf{What overall interpretation or emotional stance does the non-verbal delivery encourage?}
The non-verbal delivery encourages \textbf{moral indignation grounded in forensic certainty}. By pairing calm, courtroom-style document walkthroughs and finger-counting with incredulous pitch peaks \texttt{$\uparrow$} on the 5-minute price jump and dry scoffs at the Trust's denials, Dhruv guides the viewer to see the transaction as an open-and-shut case of institutional corruption exploiting religious faith.\par\smallskip
  \item \textbf{Whose voice or physical presence is acoustically/visually foregrounded, undermined, or silenced?}
\textbf{Documentary records (sale deeds, stamp duty receipts, witness signatures) and opposition whistleblowers} are visually foregrounded at full screen and acoustically reinforced via firm falling intonation \texttt{$\downarrow$}. \textbf{Trust General Secretary Champat Rai} is visually and acoustically undermined---his denials are introduced as \textit{"bahane"} with a smirk and followed by a \texttt{(1.5)} deadpan silence that highlights their evasiveness.\par\smallskip
  \item \textbf{What evidence supports your non-verbal reading? What complicates it (e.g., speaker baseline, format)?}
The systematic pairing of spatial hand gestures (narrow vs. wide span) and emphatic stress on time/money numerals (\texttt{5 minutes}, \texttt{2 crore}, \texttt{18.5 crore}, \texttt{[03:10]}--\texttt{[04:55]}) demonstrates how kinesics and prosody work together to construct an evidentiary indictment. What complicates this is \textbf{speaker baseline}: Dhruv Rathee uses this exact didactic cadence, finger-counting, and wide-eyed emphasis across all his explainer videos (whether on science, history, or politics), meaning his authoritative posture is part of his pedagogical brand, though the clustering of \texttt{[scoff]} vocalisations around Champat Rai is a marked evaluative shift.\par\smallskip
\end{itemize}

\bigskip\hrule\bigskip

\subsection*{Event 2. Students' protest regarding NEET paper leak (Non-Verbal \& Paralinguistic Layer)}

\begin{itemize}
  \item \textbf{Focus:} How do prosody, pauses, trailing sentences, vocalisations, and visual body language construct the legitimacy of student protesters, police use of force, and state accountability?
  \item \textbf{Outlet:} Dhruv Rathee (YouTube Channel) | \textbf{Date:} 16 June 2024
  \item \textbf{Video title:} The Dark Reality of NEET Scam | Exposed by Dhruv Rathee | \href{https://youtu.be/l-6Pm-112io}{https://youtu.be/l-6Pm-112io} | \textbf{Timestamp(s):} \texttt{[00:00]}--\texttt{[21:15]}
  \item \textbf{Analyst:} Baani Solanki (\texttt{2025114016}) | \textbf{Language:} Hinglish (Hindi in Roman script + English terms)
  \item \textbf{Format / Visual setup:} \texttt{[ ]} TV Panel Debate \texttt{[ ]} TV Bulletin \texttt{[x]} Solo YouTube Explainer \texttt{[ ]} Interview / Discussion
  \item \textbf{Primary speaker(s) observed:} Dhruv Rathee (solo studio presenter with full-screen statistical charts, FIR copies, student protest footage, and Alakh Pandey / educator interview clips)
\end{itemize}

\subsubsection*{Event-specific non-verbal questions (Event 2)}

\noindent\textbf{1.} \textbf{Prosodic \& gestural marking of participant labels ("students/bachche" vs. "protesters" vs. "upadravi/goons"): How do vocal tone, volume, facial expression, and hand gestures shift across labels? (e.g., softened, empathetic pitch for "students" vs. harsh, staccato stress, scowls, or accusatory finger-pointing on "upadraviyon", "goons", or "anti-social elements".)}\par\smallskip
\begin{itemize}
  \item \textbf{Notes:}
Dhruv exclusively uses legitimising, humanising labels---\textit{"students"}, \textit{"masoom bachche"}, \textit{"desh ka bhavishya"}, and \textit{"doctors"} (\texttt{[00:45]}, \texttt{[02:15]}, \texttt{[18:30]})---delivered with a \textbf{softened, lower-pitched empathetic register}, furrowed brow of concern, and open palms held toward his chest \texttt{\{\{gesture: open palms drawn toward chest\}\}}. Conversely, harsh staccato stress, scowls, and accusatory finger-pointing are reserved entirely for the institutional and criminal actors: \textit{"education \_MAFIA\_"}, \textit{"NTA ki \_laaparwahi\_"}, and \textit{"corrupt \_system\_"} (\texttt{[09:20]}, \texttt{[14:45]}, \texttt{[19:10]}), where he points sharply downward and strikes his palm with his fist \texttt{\{\{gesture: right fist striking open left palm\}\}}.\par\smallskip
\end{itemize}

\noindent\textbf{2.} \textbf{Delivery of students' stated demands: When the protesters' core demands (Education Minister's resignation, compensation, release of Sonam Wangchuk, exam reform) are mentioned, are they delivered slowly with deliberate pauses and finger-counting, rushed through in a flat monotone, mocked with air-quotes/scoffs, or trailed off (...)?}\par\smallskip
\begin{itemize}
  \item \textbf{Notes:}
In the concluding section (\texttt{[17:40]}--\texttt{[19:50]}), Dhruv articulates the students' core demands---(1) immediate removal/accountability of NTA leadership, (2) independent investigation into the Bihar and Gujarat paper leaks, and (3) complete transparency in raw scores and grace-mark formulas---\textbf{slowly, deliberately, and with structured finger-counting}. He pauses \texttt{(0.8)} between each point, looks directly into the lens without blinking, and delivers each demand with resolute falling intonation \texttt{$\downarrow$} (*"Ye bachchon ki koi ehsaan ki maang nahin hai (.) ye unka \_haq\_ hai$\downarrow$"\texttt{, }[18:15]`).\par\smallskip
\end{itemize}

\noindent\textbf{3.} \textbf{Where does vocal \& kinetic energy peak---grievance, street disruption, or electoral horse-race?: Track where the speaker's volume, pitch range, pace, and gesture intensity reach their climax: during accounts of the paper leak/student suicides, during street clashes and barricade breaches, or during partisan BJP-vs-Opposition sparring.}\par\smallskip
\begin{itemize}
  \item \textbf{Notes:}
Vocal and kinetic energy reaches two distinct climaxes, both centred on \textbf{student grievance and systemic betrayal}:\par\smallskip
  \item \textbf{Emotional/Sombre Climax (\texttt{[00:00]}--\texttt{[02:10]}):} Opening with the human cost of NEET preparation and student suicides, where Dhruv speaks in a hushed, slow, deeply sombre cadence (\texttt{$^\circ$ek-ek number ke liye bachche raat-raat bhar jagte hain$^\circ$}, \texttt{[01:12]}) with a heavy \texttt{[sigh]} and pained facial expression.
  \item \textbf{Indignant/Forensic Climax (\texttt{[10:15]}--\texttt{[13:50]}):} When connecting the Patna FIR (\texttt{5 May}), the Godhra cheating centre (\texttt{Rs 10 lakh} per candidate), and the Education Minister's flat denial, his volume rises sharply, his pace accelerates, and his hand gestures become rapid and forceful (\texttt{"Jab Bihar Police ki FIR mein saaf likha hai (.) toh Education Minister kis munh se keh rahe hain ki \_koi\_ leak nahin hua$\uparrow$?!"}, \texttt{[12:28]}).
\end{itemize}

\noindent\textbf{4.} \textbf{Non-verbal attribution of responsibility (sighs, stops, head shakes, pointing): When examination authorities (NTA), the Education Minister, police, opposition politicians, or students are named as responsible, what non-verbal cues accompany the attribution (exasperated sighs, mocking laughs, abrupt self-stops ---, or direct finger-pointing at the lens)?}\par\smallskip
\begin{itemize}
  \item \textbf{Notes:}
Responsibility is attributed directly to the \textbf{NTA chairperson, Education Minister Dharmendra Pradhan, and the ruling government's pattern of paper leaks (\texttt{70+} leaks in \texttt{7} years, \texttt{[15:20]})}:\par\smallskip
  \item When citing the fact that the NTA formed a committee to investigate the NTA---headed by the UPSC chair/insiders (\texttt{[08:35]})---Dhruv stops mid-sentence \texttt{---}, lets out a mocking \texttt{[chuckle]}, shakes his head, and says: \textit{"Chor hi (.) chori ki jaanch karega$\downarrow$"} (\texttt{[08:42]}) with a circular hand gesture indicating a closed loop \texttt{\{\{gesture: both index fingers rotating in a circle\}\}}.
  \item When addressing the Education Minister's initial denial (\texttt{[12:15]}), Dhruv points his index finger directly at the camera lens \texttt{\{\{gesture: sustained index-finger point at lens\}\}} with a stern, unsmiling scowl.
\end{itemize}

\noindent\textbf{5.} \textbf{Prosody of questions---open inquiry vs. loaded rhetorical challenge: Are questions about the protest delivered with level, open intonation, or with sharp rising pitch $\uparrow$, sarcastic stress, raised eyebrows, and trailing clauses (...) that presuppose protesters are violent or politically manipulated?}\par\smallskip
\begin{itemize}
  \item \textbf{Notes:}
Every single loaded rhetorical challenge in the video is aimed \textbf{upward at state authorities}, never downward at protesting students:\par\smallskip
  \item At \texttt{[05:40]}, questioning the \texttt{4 June} result timing: \textit{"Result \texttt{14} June ko aana tha, par achanak \texttt{4} June---Election Result ke din---kyon$\uparrow$ declare kiya gaya?! Kya NTA kuchh \_chhupana\_ chahti thi$\uparrow$?"} (\texttt{[05:48]}), delivered with rising pitch \texttt{$\uparrow$}, narrowed eyes, and a forward lean.
  \item By contrast, questions about the protesting students (\texttt{"Aaj jab ye bachche dhoop mein sadak par khade hain, toh inki sunne wala kaun hai$\downarrow$"}, \texttt{[16:40]}) use a falling, mournful contour \texttt{$\downarrow$} and open-palm empathy gestures that presuppose the complete legitimacy of their protest.
\end{itemize}

\noindent\textbf{6.} \textbf{Visual \& paralinguistic performance of injury and violence: How are claims of violence substantiated non-verbally on screen? (Compare camera close-ups and pained vocal tone on police bandages/fractured hands vs. visual/acoustic treatment of student head injuries, lathi-charge footage, or defensive stammers/pauses when an officer or guest is asked who used force first.)}\par\smallskip
\begin{itemize}
  \item \textbf{Notes:}
Dhruv substantiates the violence of the crisis through \textbf{close-up visuals of crying students, anguished parents, student suicide headlines, and police roughing up peaceful demonstrators} (\texttt{[01:20]}, \texttt{[16:30]}--\texttt{[17:15]}). While showing footage of police dragging and detaining protesting students (\texttt{[16:45]}), his voice drops in pitch and takes on an indignant, pained edge: \textit{"Jo bachche apna haq maangne aaye (.) unke saath \_aisa\_ salook kiya ja raha hai$\downarrow$"} (\texttt{[16:52]}), accompanied by a slow, disgusted head shake \texttt{\{\{gesture: slow head shake + furrowed brow\}\}}. Zero screen time is given to police victimhood narratives.\par\smallskip
\end{itemize}

\noindent\textbf{7.} \textbf{Asymmetric delivery of police actions vs. protester actions: Compare the vocal cadence, pauses, and hand gestures used when describing police actions (e.g., calm, matter-of-fact administrative tone on "tear gas chhode gae" / "preventive arrangement") versus protester actions (alarmist pitch, striking gestures, or ironic softening as "retaliation").}\par\smallskip
\begin{itemize}
  \item \textbf{Notes:}
Dhruv distinguishes sharply between \textbf{investigative state police (Bihar EOU)} and \textbf{street crackdown policing}:\par\smallskip
  \item \textbf{Bihar Police's EOU investigation} (\texttt{[10:20]}--\texttt{[11:40]}) is narrated with crisp, authoritative approval, pointing to the recovered burnt question-paper booklets and post-dated cheques as decisive forensic proof.
  \item \textbf{Street repression of protesters} (\texttt{[16:45]}) is narrated with moral condemnation and downward-pointing gestures, framing state force against unarmed \texttt{17}- and \texttt{18}-year-old aspirants as authoritarian bullying.
  \item \textbf{Student protesters' actions} (marching, chanting, filing petitions with educators like Alakh Pandey, \texttt{[06:40]}, \texttt{[16:30]}) are narrated with uplifting, solidaristic prosody and approving nods.
\end{itemize}

\noindent\textbf{8.} \textbf{Acoustic \& visual space granted to students' own voices: Are protesting students heard in their own voice with clear audio, steady close-ups, and empathetic listening nods, or are they reduced to noisy background chaos B-roll while studio anchors/politicians speak over them?}\par\smallskip
\begin{itemize}
  \item \textbf{Notes:}
Students and their educator-advocates (such as PhysicsWallah's Alakh Pandey and protesting NEET candidates) are given \textbf{prominent, clean, unmuted acoustic and visual space} (\texttt{[01:30]}, \texttt{[06:45]}, \texttt{[16:20]}). Their clips play at full screen with clear subtitles and zero studio cross-talk. Immediately after a clip of a distraught student saying \textit{"650 marks laane ke baad bhi college nahin milega"} (\texttt{[01:42]}), Dhruv pauses \texttt{(1.0)} in silence with a solemn, empathetic expression before speaking softly (\texttt{$^\circ$soch kar dekhiye us bachche par kya beet rahi hogi$^\circ$}, \texttt{[01:46]}).\par\smallskip
\end{itemize}

\noindent\textbf{9.} \textbf{Non-verbal shift between the NEET paper-leak issue and the street protest: Does the speaker's non-verbal demeanour shift from sympathetic concern (sombre tone, sighs, solidarity gestures) when discussing the exam leak itself to combative, authoritarian, or dismissive delivery once students march toward Parliament?}\par\smallskip
\begin{itemize}
  \item \textbf{Notes:}
\textbf{No hostile shift occurs toward the students.} Rather, as the video moves from the technical exam anomalies (\texttt{[02:30]}--\texttt{[09:00]}) to the Bihar/Gujarat paper-leak mafia (\texttt{[09:30]}--\texttt{[14:30]}) and finally to the nationwide protests (\texttt{[16:00]}--\texttt{[20:30]}), Dhruv's non-verbal demeanour escalates from \textbf{analytical concern to prosecutorial solidarity}. He explicitly frames the street protests and Supreme Court petitions as the only remaining democratic check against an unaccountable government (\texttt{[17:20]}).\par\smallskip
\end{itemize}

\subsubsection*{Apply the five non-verbal features (Event 2)}

\begin{longtable}{|p{0.31\textwidth}|p{0.31\textwidth}|p{0.31\textwidth}|}
\hline
Non-verbal feature & Timestamp + spoken anchor words + non-verbal notation & What it does in context (stance, framing, implication) \\ \hline \endhead
\textbf{1. Pauses, silences \& stops} & \texttt{[01:45]} \texttt{[(1.0) solemn silence after distraught student clip]} \textit{"$^\circ$Soch kar dekhiye (.) us bachche par kya beet rahi hogi$^\circ$"} / \texttt{[08:40]} \textit{"NTA ne jaanch ke liye committee banayi---aur uska head kaun hai? (1.0) Khud NTA$\downarrow$"} & The \texttt{(1.0)} silence after the student clip builds emotional empathy; the \texttt{(1.0)} pause before \textit{"Khud NTA$\downarrow$"} delivers a sharp satirical punchline exposing conflict of interest. \\ \hline
\textbf{2. Trailing sentences (\texttt{...})} & \texttt{[05:52]} \textit{"Chaar June---jis din poore desh ka dhyan Election Results par tha---usi din achanak NEET ka result... taaki kisi ka dhyan..."} & Trails off with a narrowed gaze and slow nod, leading the audience to complete the thought that the NTA deliberately buried the scandal under election coverage. \\ \hline
\textbf{3. Intonation, pitch \& volume} & \texttt{[04:15]} \textit{"\texttt{718} aur \texttt{719} marks---jo mathematically \_IMPOSSIBLE\_ hain$\uparrow$!"} / \texttt{[12:28]} \textit{"Kis munh se keh rahe hain ki \_KOI\_ paper leak nahin hua$\uparrow$?!"} & High-volume stress on \texttt{\_IMPOSSIBLE\_} and \texttt{\_KOI\_} combined with sharp rising pitch \texttt{$\uparrow$} transforms technical audit points into moral indictments of state dishonesty. \\ \hline
\textbf{4. Coughs, sighs \& vocalisations} & \texttt{[01:10]} \texttt{[[sigh]]} \textit{"Har saal (\texttt{24} lakh) bachche is exam mein baithte hain..."} / \texttt{[08:42]} \texttt{[[chuckle / scoff]]} \textit{"Matlab chor hi chori ki jaanch karega$\downarrow$"} & The heavy \texttt{[sigh]} establishes paternal/fraternal sorrow for students; the ironic \texttt{[chuckle / scoff]} strips the NTA's internal inquiry committee of any institutional credibility. \\ \hline
\textbf{5. Hand gestures, face \& posture} & \texttt{[08:42]} \texttt{\{\{gesture: both index fingers rotating in a closed circle\}\}} / \texttt{[12:30]} \texttt{\{\{gesture: direct index-finger point at lens + stern scowl\}\}} / \texttt{[18:10]} \texttt{\{\{gesture: sequential counting on 3 fingers\}\}} & The circular finger motion visualises institutional self-dealing; the direct lens point confronts the Education Minister; finger-counting organises student demands into a rational charter. \\ \hline
\end{longtable}

\subsubsection*{Short summary (Event 2 Non-Verbal Synthesis)}

\begin{itemize}
  \item \textbf{What overall interpretation or emotional stance does the non-verbal delivery encourage?}
The non-verbal delivery encourages \textbf{deep emotional solidarity with NEET aspirants alongside righteous civic anger at systemic state failure}. By opening with hushed, sigh-laden empathy for suicidal and exhausted students and building toward loud, finger-pointing indictment of the NTA and Education Ministry, the video mobilises the viewer to view the NEET crisis as symptomatic of broader institutional rot.\par\smallskip
  \item \textbf{Whose voice or physical presence is acoustically/visually foregrounded, undermined, or silenced?}
\textbf{Students, parents, educator-activists (Alakh Pandey), and Bihar Police FIR documents} are visually and acoustically foregrounded with full-screen clarity and respectful pauses. \textbf{Education Minister Dharmendra Pradhan and the NTA Director} are acoustically and visually undermined via scoffs, circular "self-investigation" hand gestures, and immediate juxtaposition with contradicting police evidence.\par\smallskip
  \item \textbf{What evidence supports your non-verbal reading? What complicates it (e.g., speaker baseline, format)?}
The stark acoustic contrast between the softened, low-pitch register (\texttt{$^\circ$word$^\circ$}, \textasciitilde{}115 WPM) used for student suffering (\texttt{[01:00]}--\texttt{[02:10]}) and the high-volume, staccato prosecutorial register (\textasciitilde{}150 WPM, \texttt{WORD}, \texttt{$\uparrow$}) used for the Education Minister's denials (\texttt{[12:15]}--\texttt{[13:00]}) proves that prosody is systematically deployed to assign victimhood and blame. A complicating baseline factor is that Dhruv's prosecutorial cadence, dramatic pauses, and direct camera pointing are consistent across all his political critique videos, making it a signature genre style of oppositional YouTube explainers.\par\smallskip
  \item \textbf{How do the non-verbal cues interact with the spoken words? (Do they reinforce, dramatise, soften/hedge, invert meaning via sarcasm/irony, or leave an accusation unspoken via trailing off?)}
The non-verbal cues \textbf{dramatise and reinforce} the documentary evidence (pointing at FIRs and scorecards), \textbf{invert official assurances via sarcasm} (the \texttt{[scoff]} and circular finger gesture on NTA's inquiry committee at \texttt{[08:42]}), and \textbf{leave the accusation of top-down political orchestration unspoken via trailing off} (\texttt{...}) when discussing the \texttt{4 June} election-day release (\texttt{[05:52]}) and the Gujarat/Bihar exam mafia nexus (\texttt{[14:50]}).\par\smallskip
  \item \textbf{What context, speaker baseline, or format factors support or complicate this interpretation? (e.g., Is this the speaker's normal style across all videos, or a marked shift on this specific actor/claim?)}
Compared to his Event 1 video (where he maintained a brisk, forensic cadence throughout), Event 2 displays a \textbf{markedly more emotional opening and closing baseline}, using longer silences \texttt{(1.0)} and softer vocal volume \texttt{$^\circ$...$^\circ$} to honour student trauma before shifting into his standard prosecutorial rhythm.\par\smallskip
  \item \textbf{How does another outlet perform the same moment, actor, or claim non-verbally?}
On \textbf{Republic TV} (\texttt{Event 2}), anchor Arnab Goswami also shouts at the NTA, but whenever the opposition or street protests are mentioned, Arnab pivots into frantic desk-thumping and rapid-fire interruptions \texttt{[overlap]} to attack opposition politicians (\textit{"political vultures"}). In contrast, \textbf{Aaj Tak} (\texttt{Event 2}) uses stern finger-wagging and rising pitch \texttt{$\uparrow$} to interrogate students over broken barricades and injured police. Dhruv Rathee never once turns his prosecutorial gaze against the protesters; his kinetic energy remains 100\% aligned with the students from start to finish.\par\smallskip
  \item \textbf{My interpretation, with reasons:}
Dhruv Rathee's non-verbal performance illustrates \textbf{multimodal counter-framing}: in the absence of a live studio panel, he uses prosodic contrast (hushed empathy vs. sharp prosecutorial stress), visual spatialisation (wide hand spans for price/rank jumps), and strategic reaction pauses \texttt{(1.5)} to turn official government statements into self-incriminating exhibits across both the Ram Mandir land controversy and the NEET paper leak.\par\smallskip
\end{itemize}

\end{document}
